%% kapa: exam=1 round=27 year=2016 period=1 form=A subject=civil_law subject_ko=민법 seq=13 number=13
%% source: Q-Net 감정평가사 기출문제 (cst00309) artlSeq=5169402 file=extracted/1008/5169402/제27회 감정평가사 1차 문제/제27회 감정평가사 1교시문제(A형).pdf
%% conversion: kordoc 4.24.51 -> tools/gen1.py (auto); review=proofread-sonnet; flags=none
\begin{question}{13}
복대리에 관한 설명으로 옳지 않은 것은?
\begin{choices}
  \item 임의대리인은 본인의 승낙이나 부득이한 사유가 없으면, 복대리인을 선임하지 못한다.
  \item 복대리인은 제3자에 대하여도 대리인과 동일한 권리의무가 있다.
  \item 임의대리인은 본인의 지명에 의해서도 복대리인을 선임할 수 있다.
  \item 대리인의 대리권이 소멸하면 복대리인의 대리권도 소멸한다.
  \item 복대리인은 대리인이 본인의 명의로 선임한 본인의 대리인이다.
\end{choices}
\end{question}
\answer{5}
